\documentclass{situationist-preprint}

\title{Successors and Survival}
\msid{PID-14-b}
\msfamily{PID-14}
\msdiscipline{Personal Identity and the Self}
\mstype{Companion manuscript}
\mskeywords{mind uploading; successors; what matters; prudential concern; Parfit}

\begin{document}
\maketitle

\begin{abstract}
PID-14-a argues that uploading creates a psychologically continuous successor, not the original person, and that the successor has much of what matters in survival. We ask what the original has reason to want for her successor. We argue that she has reason to want for it nearly what she wants for her own near future: that its life go well, that it carry on her projects and relationships, and that it be free to change. Because identity carries little weight, her reasons differ from self-interested ones mainly in two respects: she has reason to want the successor's standing as a distinct person to be recognized, and she has no reason to prefer branching to its absence for the sake of the successor.
\end{abstract}

\section{The successor's claim on the original}

By PID-14-a, Theorems~3.1 and~3.2, an upload is not the original person but, if thinking is functional (MND-08), a person psychologically continuous with her: a successor (Definition~2.2). Discussions of uploading have mostly asked whether it is survival \citep{Chalmers2010}; once that is answered in the negative, the question becomes what the original should want for the one who comes after her. At the time of upload it is as strongly connected with her as she is with herself the next morning. By PID-06-a, connectedness carries most of what matters, and identity at most a small weight $\varepsilon$ (Theorem~3.3, Corollary~3.4).

\begin{proposition}\label{p:want}
The original has prudential reasons, nearly as strong as those she has regarding her own near future, to want her successor's life to go well.
\end{proposition}

\begin{proof}
What matters is the strength of the prudential reasons a relation gives a person for concern about a future being (PID-06-a, Definition~2.4). For the successor, identity is absent, continuity holds and connectedness is near its maximum, so what matters is at least $1-\varepsilon$, against $1$ for her near future (PID-06-a, Principle~P1, which applies because the successor's relation to her is, in its connections, the relation of a fission product to its original).
\end{proof}

What it is for the successor's life to go well, from the original's point of view, is shaped by her values, which the successor shares. She has reason to want it to carry on her projects and care for those she cares for, as she would want this of her own future. The reasons decline as the successor diverges from her, since connectedness declines (PID-06-b); so they bear most strongly on the successor's near future.

\section{Where the reasons differ}

\paragraph{Standing.} Since the successor is not she, it is a person with its own standing, not her continuation under another name. She has reason to want this recognized: that the successor not be held to her promises as if it had made them, that it be free to revise the projects it inherits, and that those who knew her not treat it as a substitute. Her concern for it, though prudential in strength, is concern for another, and respect for another's standing is part of it. The successor's own obligations from her acts are governed, like those of an adult toward her childhood self (PID-07), by connectedness and not by identity, which it lacks: strong at first, weakening as it diverges.

\paragraph{Branching.} In nondestructive uploading she continues alongside her successor, and the process branches (PID-14-a, Proposition~3.4). For the successor's sake, she has no reason to prefer this to destructive uploading, since what the successor has does not depend on whether she exists (PID-03-a, Assumption~2.6). She may have reasons against it for the sake of goods of uniqueness, such as her partner's having one spouse, which carry the small residual weight of identity. She also has a reason for it: in nondestructive uploading she keeps her identity and her life, and the upload is a successor in either case. Destructive uploading is therefore not something she has reason to choose for the successor's benefit.

\begin{objection}
If the successor has nearly all of what matters, the original should be indifferent between continuing and being replaced by it.
\end{objection}

\begin{reply}
Nearly indifferent, with respect to what matters in survival, and only in respect of the near future, since the successor's connectedness with her declines as it diverges, as does hers with her own future self. She is not indifferent in other respects: her own continuation preserves identity and the goods of uniqueness, and avoids the risks of a technology whose fidelity cannot be checked from within. \citet[ch.~12]{Parfit1984} drew a similar conclusion about teletransportation; the difference here is that the upload is never the original, whatever the timing.
\end{reply}

\section{Conclusion}

The original has reason to want for her successor nearly what she wants for her own near future, with two differences that follow from its not being she: she has reason to want it recognized as a person in its own right, and no reason, for its sake, to choose her own destruction.

\begin{thebibliography}{99}\small
\bibitem[Parfit(1984)]{Parfit1984} Parfit, D. (1984). \emph{Reasons and Persons}. Oxford University Press.
\bibitem[Chalmers(2010)]{Chalmers2010} Chalmers, D.~J. (2010). The singularity: A philosophical analysis. \emph{Journal of Consciousness Studies}, 17(9--10), 7--65.
\end{thebibliography}

\end{document}
